% META: {"id":"1SPE-GEOREP-CO-021","chapitre":"1SPE-GEOMETRIE-REPEREE","type_objet":"corrige","exercice_ref":"1SPE-GEOREP-EX-021","capacites_codes":["C3"],"sources_inspiration":[],"status":"approved","origin":"generated","status_history":["generated","audited","accepted","approved"],"release_acceptance":"RELEASE_OWNER_BATCH_ACCEPTANCE","acceptance_closure_digest":"sha256:9b3ccf9a81c5520fbb7e03b7b2d3e7bf2057a02834b6d908bf3be5f49f3b553f","acceptance_receipt_digest":"sha256:4fad86f0282e0fa4e0419cca1ad6379ae7af5a280c04a4a90880f90a1c044242"}
% BEGIN-VERIFY
% from sympy import *
% x, y = symbols('x y')
% # Cercle centre (2,-3), rayon 5
% # (x-2)^2 + (y+3)^2 = 25
% assert (0-2)**2 + (0+3)**2 == 13  # (0,0) pas sur le cercle
% assert (2-2)**2 + (2+3)**2 == 25  # (2,2) sur le cercle
% END-VERIFY
\begin{corrige}{1SPE-GEOREP-EX-021}

\textbf{Question 1.} $(x - 2)^2 + (y + 3)^2 = 25$.

\textbf{Question 2.} $(2 - 2)^2 + (2 + 3)^2 = 0 + 25 = 25$. Donc $A \in \mathcal{C}$.

\end{corrige}
